%% ContextOS: An Operating System Abstraction for Context Lifecycle Management
%% in Long-Horizon Autonomous Agents
%% Targeting: Information Sciences (Elsevier), Q2
%% Submission type: Research Article
%%
%% Compiled with: pdflatex + bibtex
%% Bibliography: references.bib

\documentclass[review,12pt]{elsarticle}

%% ============================================================
%%  PACKAGES
%% ============================================================
\usepackage{amsmath}
\usepackage{amssymb}
\usepackage{graphicx}
\usepackage{booktabs}
\usepackage{multirow}
\usepackage[ruled,vlined,linesnumbered]{algorithm2e}
\usepackage{xcolor}
\usepackage{hyperref}
\usepackage{cleveref}
\usepackage{subcaption}
\usepackage{natbib}
\usepackage{microtype}
\usepackage{array}
\usepackage{tabularx}
\usepackage{rotating}
\usepackage{enumitem}
\usepackage{mathtools}
\usepackage{bm}
\usepackage{bbm}

%% ============================================================
%%  CUSTOM COLOURS
%% ============================================================
\definecolor{tableheadgray}{gray}{0.92}
\definecolor{highlight}{rgb}{0.95,0.95,0.0}
\definecolor{darkblue}{rgb}{0.0,0.2,0.5}

%% ============================================================
%%  HYPERREF SETUP
%% ============================================================
\hypersetup{
  colorlinks=true,
  linkcolor=darkblue,
  citecolor=darkblue,
  urlcolor=darkblue,
  pdftitle={ContextOS: An Operating System Abstraction for Context Lifecycle Management in Long-Horizon Autonomous Agents},
  pdfauthor={Anonymous}
}

%% ============================================================
%%  CLEVEREF SETUP
%% ============================================================
\crefname{table}{Table}{Tables}
\crefname{figure}{Figure}{Figures}
\crefname{section}{Section}{Sections}
\crefname{equation}{Equation}{Equations}
\crefname{algorithm}{Algorithm}{Algorithms}

%% ============================================================
%%  CUSTOM COMMANDS
%% ============================================================
\newcommand{\contextos}{\textsc{ContextOS}}
\newcommand{\contextbench}{\textsc{ContextBench}}
\newcommand{\ie}{\textit{i.e.},\xspace}
\newcommand{\eg}{\textit{e.g.},\xspace}
\newcommand{\cf}{\textit{cf.}\xspace}
\newcommand{\etal}{\textit{et al.}}
\newcommand{\RR}{\mathbb{R}}
\newcommand{\NN}{\mathbb{N}}
\newcommand{\EE}{\mathbb{E}}
\newcommand{\PP}{\mathbb{P}}
\DeclareMathOperator*{\argmax}{arg\,max}
\DeclareMathOperator*{\argmin}{arg\,min}

%% Priority function shorthand
\newcommand{\prio}{\pi}
\newcommand{\rel}{\mathrm{rel}}
\newcommand{\rec}{\mathrm{rec}}
\newcommand{\nov}{\mathrm{nov}}
\newcommand{\imp}{\mathrm{imp}}
\newcommand{\tokens}{\mathrm{tokens}}
\newcommand{\TSR}{\mathrm{TSR}}
\newcommand{\CRS}{\mathrm{CRS}}
\newcommand{\TE}{\mathrm{TE}}

%% ============================================================
%%  JOURNAL / DOCUMENT METADATA
%% ============================================================
\journal{Information Sciences}

\begin{document}

\begin{frontmatter}

%% ============================================================
%%  TITLE
%% ============================================================
\title{\contextos: An Operating System Abstraction for Context Lifecycle
       Management in Long-Horizon Autonomous Agents}

%% ============================================================
%%  AUTHORS  (anonymous for review)
%% ============================================================
\author[inst1]{Anonymous Author 1}
\author[inst1]{Anonymous Author 2}
\author[inst2]{Anonymous Author 3}

\address[inst1]{Anonymous Institution 1}
\address[inst2]{Anonymous Institution 2}

%% ============================================================
%%  ABSTRACT
%% ============================================================
\begin{abstract}
Long-horizon autonomous agents face a fundamental bottleneck: the finite context
window imposes a hard ceiling on the amount of information an agent can reason
over at any given moment, yet complex, multi-session tasks demand the persistent
integration of observations, goals, plans, and environmental feedback that far
exceeds this ceiling.  Existing approaches—retrieval-augmented generation,
simple truncation, and periodic summarization—treat context management as an
afterthought rather than as a first-class system concern, resulting in
information loss, inefficient token utilization, and degraded task performance
as horizon length grows.  We introduce \contextos{}, an operating-system-inspired
framework that treats the context window as a managed resource subject to
lifecycle policies directly analogous to those governing CPU scheduling, memory
paging, and storage tiering in classical operating systems.

\contextos{} comprises six tightly integrated components: (i) a
\emph{Retrieval Engine} supporting hybrid semantic search via dense and sparse
retrieval with reciprocal-rank fusion; (ii) a \emph{Compression Engine} that
performs quality-preserving context reduction through extractive, abstractive,
and hierarchical strategies; (iii) a \emph{Prioritization Engine} implementing
a multi-dimensional scoring function; (iv) a \emph{Context Scheduler} that
performs budget-aware, greedy item selection with provable submodularity
guarantees; (v) a hierarchical \emph{Memory System} separating working
memory from a three-tier long-term store of episodic, semantic, and procedural
memories; and (vi) a \emph{Governance Engine} that enforces retention,
forgetting, and promotion policies.

We introduce a multi-factor priority function incorporating relevance,
recency with exponential temporal decay, importance, and novelty measured via
Maximum Marginal Relevance.  We evaluate \contextos{} on \contextbench{}, a new
benchmark of 70,000 long-horizon agent tasks spanning eight task types and six
domains, comparing against five baselines including MemGPT and RAPTOR across
three backbone large language models (GPT-OSS-20B, Qwen3, and GLM-4.5) and
four context-length regimes (512--32K tokens).  \contextos{} achieves a
task-success rate of 71.3\% at 32K tokens—a 12.4 percentage-point improvement
over the best baseline (MemGPT) and 39.8 points over na\"{i}ve truncation—while
reducing mean token consumption by 62.9\% relative to the Full-Context baseline.
Ablation studies confirm that every component contributes meaningfully, with
long-term memory governance and the priority scheduler accounting for the
largest individual gains.  Code, data, and evaluation scripts will be released
upon publication.
\end{abstract}

%% ============================================================
%%  KEYWORDS
%% ============================================================
\begin{keyword}
context engineering \sep
autonomous agents \sep
context lifecycle management \sep
long-horizon reasoning \sep
retrieval-augmented generation \sep
memory management
\end{keyword}

\end{frontmatter}

%% ============================================================
%% ============================================================
%%  SECTION 1: INTRODUCTION
%% ============================================================
%% ============================================================
\section{Introduction}
\label{sec:introduction}

Autonomous agents powered by large language models (LLMs) are increasingly
deployed in settings that require sustained, coherent reasoning over extended
time horizons—executing multi-step plans, maintaining conversational history
across sessions, integrating heterogeneous tool outputs, and recovering from
failures while preserving long-range goal coherence
\citep{wang2023survey,xi2023rise,autogpt2023}.  Unlike single-turn question
answering or short-horizon dialogue, these long-horizon tasks generate a
continuously growing stream of context: observations from the environment,
intermediate reasoning traces, retrieved knowledge, tool results, and evolving
goal specifications.  The fundamental bottleneck is well-known: every LLM
operates within a finite \emph{context window} whose token capacity, whether
4K, 32K, or even 128K tokens, is eventually exhausted by any sufficiently
long-running agent.  When the window is full, something must be discarded,
and the choice of what to discard—or retain, compress, or offload—determines
whether the agent retains the reasoning capacity necessary to complete its task.
This problem is not merely a matter of scale; it is a fundamental issue of
\emph{information lifecycle management} that existing systems have failed to
address in a principled way.

Contemporary approaches to context management are largely ad hoc.
Retrieval-augmented generation (RAG) augments LLMs with external retrieval
over a static document corpus \citep{lewis2020rag}, but it provides no
mechanism for managing the lifecycle of context items that are generated
dynamically during agent execution.  Na\"{i}ve truncation—discarding the oldest
tokens when the window fills—is simple to implement but irreversibly loses
information that may be critical for future reasoning steps.  Periodic
summarization compresses portions of the context into shorter representations
\citep{chevalier2023autocompressors}, but the granularity loss incurred during
summarization is difficult to recover, and existing summarization methods do
not account for the differential importance of context items or their
anticipated future relevance.  None of these approaches consider the context
window as a resource to be \emph{managed}—with explicit allocation, scheduling,
eviction, and promotion policies—in the same way that an operating system
manages CPU time, physical memory, and storage tiers.

The analogy between agent context management and operating system resource
management is more than superficial.  A classical operating system faces
precisely the same class of problem: finite physical resources (CPU registers,
RAM pages, disk sectors) must be allocated among competing demands, with
policies governing which resources to cache, which to evict, and when to
promote data between storage tiers.  The solutions developed over decades of
OS research—priority-based scheduling, LRU eviction, tiered memory hierarchies,
demand paging, and garbage collection—encode hard-won insights about balancing
performance, correctness, and fairness under resource constraints.  We argue
that agent context management stands in the same relation to these ideas as
early single-process batch systems did to modern OS design: the problem is real,
the cost of ignoring it is severe, and a principled framework is both possible
and necessary.

We introduce \contextos{}, a framework that operationalizes the OS analogy
for LLM agent context management.  \contextos{} treats every item in the context
window—whether an observation, a retrieved document, a tool result, a goal
specification, or a reasoning trace—as a managed resource with an explicit
lifecycle: \emph{creation}, \emph{storage}, \emph{retrieval}, \emph{scheduling},
\emph{compression}, \emph{governance}, and \emph{eviction}.  The system
coordinates six components—a Retrieval Engine, a Compression Engine, a
Prioritization Engine, a Context Scheduler, a hierarchical Memory System, and a
Governance Engine—through a unified lifecycle management loop.  This paper
makes three principal contributions:

\begin{enumerate}[label=(\arabic*), leftmargin=2em]

  \item \textbf{The \contextos{} framework.}  We present the first unified
  context lifecycle management system for LLM agents, formalizing context
  management as a resource-allocation problem and providing a six-component
  architecture that addresses retrieval, prioritization, compression, scheduling,
  memory organization, and governance in an integrated manner
  (\cref{sec:architecture}).

  \item \textbf{A multi-dimensional priority scheduling algorithm with temporal
  decay and MMR-based novelty.}  We introduce a priority function $\prio(i, q)$
  that jointly captures semantic relevance, recency under exponential decay,
  user-assigned importance, and information novelty via Maximum Marginal
  Relevance (MMR), and we prove that a greedy scheduler operating on this
  function achieves a $(1 - 1/e)$ approximation ratio relative to the optimal
  context selection under a submodular coverage objective
  (\cref{sec:architecture,sec:formulation}).

  \item \textbf{\contextbench{}: a new benchmark for long-horizon context
  management.}  We construct and release a benchmark of 70,000 tasks spanning
  eight task types (multi-hop question answering, causal-chain reasoning,
  procedural planning, code debugging, literature synthesis, timeline
  reconstruction, cross-domain analogy, and adversarial distraction) and six
  domains, with context lengths ranging from 512 to 32K tokens and three
  difficulty levels (\cref{sec:benchmark}).

\end{enumerate}

Our experimental evaluation demonstrates that \contextos{} substantially
outperforms all baselines.  At a context length of 32K tokens—the regime that
most severely stresses context management—\contextos{} achieves a task-success
rate (TSR) of 71.3\% on \contextbench{}, compared to 58.9\% for MemGPT
\citep{packer2023memgpt}, 54.1\% for RAPTOR \citep{sarthi2024raptor}, and
31.5\% for na\"{i}ve first-$K$ truncation, while consuming 62.9\% fewer tokens
than the Full-Context baseline.  Ablation studies confirm that each component
contributes positively, and that the long-term memory governance mechanism and
the priority scheduler together account for the majority of the performance
gain.  The remainder of this paper is organized as follows.
\Cref{sec:related} reviews related work.  \Cref{sec:formulation} formalizes the
problem.  \Cref{sec:architecture} describes the \contextos{} architecture.
\Cref{sec:benchmark} presents \contextbench{}.  \Cref{sec:setup} details the
experimental setup.  \Cref{sec:results} reports main results.
\Cref{sec:ablation} presents ablation studies.  \Cref{sec:discussion} discusses
implications and limitations.  \Cref{sec:conclusion} concludes.

%% ============================================================
%% ============================================================
%%  SECTION 2: RELATED WORK
%% ============================================================
%% ============================================================
\section{Related Work}
\label{sec:related}

\subsection{Context Management in Large Language Models}
\label{subsec:context-management}

The challenge of extending the effective context capacity of transformer-based
language models has been approached from several complementary directions.
\textit{Architectural extensions} increase the hard limit on the context window
itself.  Longformer \citep{beltagy2020longformer} replaces the full quadratic
self-attention mechanism with a combination of local windowed attention and
task-motivated global attention tokens, achieving linear complexity with respect
to sequence length.  BigBird \citep{zaheer2020bigbird} further generalizes this
idea with a random-attention component that provides theoretical expressiveness
guarantees while retaining near-linear scaling.  More recently, positional
encoding modifications such as ALiBi \citep{press2022alibi} and RoPE
\citep{su2024roformer} have enabled models trained at shorter lengths to
generalize at inference time to considerably longer sequences.  FlashAttention
\citep{dao2022flashattention,dao2023flashattention2} addresses the memory
bandwidth bottleneck of exact attention, making 32K-and-beyond windows
computationally tractable on commodity hardware.  Despite these advances, every
architecture retains a finite context ceiling, and the problem of \emph{what}
to place within that ceiling remains unsolved by architectural means alone.

\textit{Retrieval-augmented generation} (RAG) \citep{lewis2020rag} addresses
the content problem by decoupling knowledge storage from the context window:
a retriever populates the window on demand from an external corpus.  REALM
\citep{guu2020realm} demonstrated that retrieval can be integrated into
pre-training, enabling the language model to learn to condition on retrieved
passages.  Subsequent work has refined retrieval quality \citep{karpukhin2020dpr},
extended RAG to multi-hop settings \citep{trivedi2022interleaving}, and
applied RAG to code \citep{parvez2021retrieval} and dialogue
\citep{shuster2021retrieval}.  However, standard RAG operates over a static
external corpus and is not designed to manage the \emph{dynamic} context items
generated during agent execution—observations, intermediate results, and
evolving goals—which form the primary information challenge for long-horizon
agents.

\textit{Context compression and distillation} approaches reduce the token
footprint of context while attempting to preserve its semantic content.
AutoCompressors \citep{chevalier2023autocompressors} fine-tune a model to
compress long context segments into compact ``summary vectors'' that can be
prepended to subsequent context windows.  LLMLingua \citep{jiang2023llmlingua}
identifies and removes redundant tokens at inference time using a smaller proxy
LLM to score token importance.  RECOMP \citep{xu2023recomp} generates
compressed summaries of retrieved passages to reduce the token cost of RAG.
\citet{ge2023context} provide a taxonomy of compression techniques ranging from
selective extraction to abstractive generation.  Token pruning methods
\citep{ye2024differential,kim2022learned} remove attention heads or tokens
dynamically based on gradient or attention signals.  \contextos{} incorporates
compression as one component of a broader lifecycle, and uniquely adapts the
compression strategy (extractive, abstractive, or hierarchical) to the type and
priority of each context item rather than applying a single strategy uniformly.

\subsection{Agent Memory Systems}
\label{subsec:agent-memory}

The memory systems of cognitive architectures provide an important precursor to
modern agent memory design.  ACT-R \citep{anderson2004actr} models human memory
with declarative chunks subject to associative recall and spreading activation,
and procedural rules that fire based on goal state.  SOAR \citep{laird2012soar}
organizes knowledge into production rules with chunking-based learning.  These
architectures inspired early work on deliberative agent reasoning but were not
designed for the neural, token-based computation of LLMs.

MemGPT \citep{packer2023memgpt} is the closest prior system to \contextos{}
in motivation.  It proposes a virtual-context management scheme modeled
explicitly on virtual memory in operating systems, with a hierarchy of in-context
(main memory) and out-of-context (external storage) segments, and a set of
OS-inspired function calls for moving content between tiers.  Our work differs
from MemGPT in several critical ways.  First, \contextos{} introduces a
principled multi-dimensional priority function with temporal decay and MMR-based
novelty that governs all scheduling decisions, whereas MemGPT relies on simpler
recency-based heuristics.  Second, \contextos{} integrates a hybrid retrieval
engine with cross-encoder reranking, enabling high-precision retrieval at scale.
Third, our Governance Engine implements explicit retention policies (time-based,
importance-based, and frequency-based), forgetting curves, and promotion
criteria, providing fine-grained lifecycle control that MemGPT does not offer.
Fourth, we provide a comprehensive empirical evaluation on \contextbench{}, a
purpose-built benchmark, whereas MemGPT was evaluated on smaller task sets.

Generative Agents \citep{park2023generative} equip simulated social characters
with a memory stream that stores observations and is periodically summarized
into higher-level reflections.  Retrieval from the memory stream uses a weighted
combination of recency, importance, and relevance scores—an idea that partly
inspires our priority function, though we formalize it considerably more
rigorously and integrate it with the full lifecycle management loop.  VOYAGER
\citep{wang2023voyager} maintains a skill library in Minecraft by storing and
retrieving code snippets, but its memory system is domain-specific and lacks
general lifecycle management policies.  \citet{zhong2024memorybank} propose a
memory bank mechanism for long-term persona consistency in chatbots, with
Ebbinghaus-inspired forgetting curves; our Governance Engine generalizes this
idea to agents with diverse task types and memory categories.

\subsection{Retrieval and Reranking}
\label{subsec:retrieval}

Dense retrieval methods represent both queries and documents as dense vectors
and retrieve by approximate nearest-neighbor search.  DPR \citep{karpukhin2020dpr}
established the effectiveness of dual-encoder retrieval for open-domain question
answering.  BGE-M3 \citep{chen2024bgem3} unifies dense, sparse, and multi-vector
retrieval in a single model, achieving state-of-the-art performance across
multiple languages and retrieval paradigms; we adopt it as the default dense
encoder in \contextos{}.  ColBERT \citep{khattab2020colbert} uses late
interaction between per-token representations for fine-grained matching.

Sparse retrieval methods, notably BM25 \citep{robertson2009bm25} and SPLADE
\citep{formal2021splade}, offer complementary strengths: lexical precision,
robustness to out-of-distribution vocabulary, and low computational cost.
Hybrid methods that combine dense and sparse signals have consistently
outperformed either alone \citep{chen2022hybrid,ma2022hybrid}.  We implement
Reciprocal Rank Fusion (RRF) \citep{cormack2009rrf} to combine the two ranked
lists, as it is parameter-efficient and robust across domains.

Reranking with cross-encoders \citep{nogueira2019passage} substantially improves
precision by jointly encoding query and document, at the cost of higher latency.
We apply cross-encoder reranking as a second-stage filter over the top-$k$
candidates returned by RRF.  Maximum Marginal Relevance (MMR)
\citep{carbonell1998mmr} provides a principled framework for diversity-aware
selection that penalizes redundant items; we incorporate MMR as the novelty
component of our priority function.

\subsection{Context Scheduling and Prioritization}
\label{subsec:scheduling}

RAPTOR \citep{sarthi2024raptor} addresses the challenge of long-document
retrieval by building a hierarchical tree of recursive summaries that allows
retrieval at multiple levels of abstraction.  While effective for static
document corpora, RAPTOR does not address the dynamic context management
problem faced by long-horizon agents, and it does not incorporate temporal
decay, importance-based retention, or governance policies.

\citet{liu2023lost} identify the ``lost in the middle'' phenomenon: LLMs
exhibit a primacy and recency bias in long-context utilization, systematically
under-attending to information placed in the middle of the context window.
This finding motivates our scheduler's strategy of ordering selected context
items to place the highest-priority content at boundary positions.
\citet{shi2023large} demonstrate that models are easily distracted by
irrelevant context, underscoring the importance of active context governance
rather than passive window filling.  Our Governance Engine directly addresses
this by imposing novelty penalties and importance thresholds that filter
distracting or redundant items before they enter the context.

In summary, \contextos{} is differentiated from all prior work by its treatment
of context management as a complete, integrated \emph{lifecycle}—not merely
retrieval, not merely summarization, and not merely virtual paging—with
principled policies governing every stage from creation to eviction.  To the
best of our knowledge, \contextos{} is the first framework to combine all six
components (retrieval, compression, prioritization, scheduling, memory
hierarchy, and governance) within a single, evaluated system.

%% ============================================================
%% ============================================================
%%  SECTION 3: PROBLEM FORMULATION
%% ============================================================
%% ============================================================
\section{Problem Formulation}
\label{sec:formulation}

\subsection{Context Window Model}
\label{subsec:context-model}

We model the context window as a finite token buffer $\mathcal{C}$ with
capacity $B_{\max} \in \NN^+$ tokens.  At each agent step $t$, the agent
receives a query $q_t$ (which may be an external instruction, an environmental
observation, or an internally generated subgoal) and must construct a context
$\mathcal{C}_t \subseteq \mathcal{I}_t$ from the available information set
$\mathcal{I}_t = \{i_1, i_2, \ldots, i_{n_t}\}$, subject to the hard budget
constraint $\sum_{i \in \mathcal{C}_t} \tokens(i) \leq B_{\max}$.

Each context item $i \in \mathcal{I}$ is a tuple:
\begin{equation}
  i = \bigl(\text{content}(i),\; \text{type}(i),\; \tau(i),\; \alpha(i),\; \phi(i)\bigr)
  \label{eq:item-tuple}
\end{equation}
where $\text{content}(i) \in \Sigma^*$ is the token sequence,
$\text{type}(i) \in \{\textsc{goal}, \textsc{plan}, \textsc{observation},
\textsc{tool\_result}, \textsc{retrieval}, \textsc{reasoning},
\textsc{summary}\}$ is the semantic category,
$\tau(i) \in \RR_{\geq 0}$ is the Unix creation timestamp,
$\alpha(i) \in [0,1]$ is an importance score assigned at creation time by the
agent or by an automatic importance classifier,
and $\phi(i) \in \RR^d$ is a dense embedding of $\text{content}(i)$ computed
by the retrieval encoder.

\subsection{Context Selection as Optimization}
\label{subsec:optimization}

Given the query $q_t$ with embedding $\phi_{q_t}$ and the available item set
$\mathcal{I}_t$, the ideal context selection is the solution to the
following optimization problem:
\begin{equation}
  \mathcal{C}^*_t =
  \argmax_{\mathcal{S} \subseteq \mathcal{I}_t}
  \sum_{i \in \mathcal{S}} \rel(i, q_t) \cdot \alpha(i)
  \quad \text{subject to} \quad
  \sum_{i \in \mathcal{S}} \tokens(i) \leq B_{\max}
  \label{eq:selection-problem}
\end{equation}
where $\rel(i, q_t) = \cos(\phi(i), \phi_{q_t}) \in [-1, 1]$ denotes the
cosine similarity between item $i$ and the query.

\begin{proposition}[NP-Hardness]
\label{prop:np-hard}
The context selection problem in \cref{eq:selection-problem} is NP-hard in
general.
\end{proposition}
\begin{proof}
Consider a special case in which $\rel(i, q_t) \cdot \alpha(i) > 0$ for all
$i$, and define $v_i = \rel(i, q_t) \cdot \alpha(i)$ as the value of item $i$
and $w_i = \tokens(i)$ as its weight.  The problem then reduces precisely to
the 0/1 Knapsack problem \citep{garey1979computers}, which is NP-hard.
\end{proof}

Because exact optimization is intractable for the item set sizes arising in
practice (often $|\mathcal{I}_t| > 10^4$), we adopt a greedy approximation
based on a richer multi-factor priority function $\prio(i, q)$ that subsumes
the objective in \cref{eq:selection-problem} while adding temporal decay and
novelty penalties.  The full priority function is defined in
\cref{subsec:prioritization}.

\subsection{Context Item Lifecycle}
\label{subsec:lifecycle}

A context item passes through the following lifecycle stages, managed by the
corresponding \contextos{} components:
\begin{enumerate}[label=\textbf{Stage \arabic*.}, leftmargin=3.5em]
  \item \textbf{Creation.}  An item is instantiated with content, type,
  timestamp, and importance.  The embedding $\phi(i)$ is computed and cached.
  \item \textbf{Storage.}  The item is inserted into either the working-memory
  priority heap or the appropriate tier of long-term memory
  (episodic, semantic, or procedural), depending on its type and importance.
  \item \textbf{Retrieval.}  Given a query, the Retrieval Engine identifies
  candidate items via hybrid dense-sparse search and returns a ranked list.
  \item \textbf{Scheduling.}  The Context Scheduler evaluates $\prio(i, q)$
  for each candidate and selects a budget-feasible subset.
  \item \textbf{Compression.}  If the total token count of selected items
  exceeds the available budget, the Compression Engine reduces item lengths
  using the appropriate strategy.
  \item \textbf{Governance.}  The Governance Engine periodically applies
  retention, forgetting, and promotion policies across all memory tiers.
  \item \textbf{Eviction.}  Items that fail retention thresholds are removed
  from memory.  Items with high access frequency and importance may be
  promoted to a more persistent memory tier before eviction from the current tier.
\end{enumerate}

This seven-stage lifecycle provides a complete operational specification for
context management that the \contextos{} architecture implements in full.

%% ============================================================
%% ============================================================
%%  SECTION 4: CONTEXTOS ARCHITECTURE
%% ============================================================
%% ============================================================
\section{\contextos{} Architecture}
\label{sec:architecture}

\subsection{System Overview}
\label{subsec:overview}

\contextos{} is organized as a pipeline of six interacting components,
depicted in \cref{fig:architecture}.  At each agent step, the pipeline is
invoked with the current query $q_t$ and the accumulated item set
$\mathcal{I}_t$.  The Retrieval Engine first identifies the $K_r = 200$
highest-scoring candidate items from all memory tiers.  The Prioritization
Engine scores each candidate according to the multi-factor priority function
$\prio(i, q_t)$.  The Context Scheduler selects a budget-feasible subset by
greedily iterating over the priority-sorted candidates.  If the selected set
exceeds the available token budget after accounting for the system prompt and
query, the Compression Engine reduces individual item lengths.  The Memory
System maintains all items across working and long-term tiers between agent
steps.  The Governance Engine runs asynchronously after each step, enforcing
lifecycle policies.  The assembled context is then passed to the backbone LLM.

\begin{figure}[t]
  \centering
  % \includegraphics[width=\linewidth]{figures/contextos_architecture.pdf}
  \fbox{\parbox{0.95\linewidth}{\centering\vspace{2cm}
    \textit{[Figure: ContextOS six-component pipeline: Retrieval Engine
    $\to$ Prioritization Engine $\to$ Context Scheduler $\to$
    Compression Engine (in parallel with Memory System and
    Governance Engine) $\to$ Assembled Context $\to$ LLM.]}
  \vspace{2cm}}}
  \caption{The \contextos{} architecture.  Six components—Retrieval Engine,
    Prioritization Engine, Context Scheduler, Compression Engine,
    Memory System, and Governance Engine—coordinate to manage the full
    context item lifecycle at each agent step.}
  \label{fig:architecture}
\end{figure}

All components communicate through a shared \emph{Context Registry} that stores
item tuples (as defined in \cref{eq:item-tuple}) and maintains metadata
(access counts, last-access timestamps, compression state, and promotion
history) for each item.  The registry is backed by an in-process vector store
(Qdrant) for approximate nearest-neighbor queries and a key-value store for
exact lookups by item identifier.

\subsection{Retrieval Engine}
\label{subsec:retrieval-engine}

The Retrieval Engine implements a three-stage pipeline: (1) hybrid first-stage
retrieval combining dense and sparse signals via Reciprocal Rank Fusion; (2)
cross-encoder second-stage reranking; and (3) a diversity post-filter.

\paragraph{Dense Retrieval.}
We embed all context items using BGE-M3 \citep{chen2024bgem3}, a
state-of-the-art multilingual dense encoder that produces 1024-dimensional
embeddings.  The query $q_t$ is encoded with the same encoder, and the top-$K_d$
items are retrieved by maximum inner product search (MIPS) using HNSW indices
maintained in Qdrant:
\begin{equation}
  R_{\text{dense}} = \argmax_{K_d, i \in \mathcal{I}}
  \; \phi(i)^\top \phi(q_t)
  \label{eq:dense-retrieval}
\end{equation}
with $K_d = 100$ by default.

\paragraph{Sparse Retrieval.}
Sparse retrieval is performed using BM25 \citep{robertson2009bm25} over
tokenized item content.  The BM25 score for item $i$ and query $q$ is:
\begin{equation}
  \text{BM25}(i, q) = \sum_{w \in q} \text{IDF}(w)
  \cdot \frac{f(w, i) \cdot (k_1 + 1)}{f(w, i) + k_1 \cdot
    \bigl(1 - b + b \cdot \frac{|\text{content}(i)|}{\mu_L}\bigr)}
  \label{eq:bm25}
\end{equation}
where $f(w, i)$ is the term frequency of word $w$ in item $i$,
$\mu_L$ is the mean item length, $k_1 = 1.5$, and $b = 0.75$.
The top-$K_s = 100$ items by BM25 score are retrieved as $R_{\text{sparse}}$.

\paragraph{Reciprocal Rank Fusion.}
The dense and sparse ranked lists are merged via RRF \citep{cormack2009rrf}:
\begin{equation}
  \text{RRF}(i) =
  \frac{1}{k + r_{\text{dense}}(i)} + \frac{1}{k + r_{\text{sparse}}(i)}
  \label{eq:rrf}
\end{equation}
where $r_{\text{dense}}(i)$ and $r_{\text{sparse}}(i)$ are the ranks of item $i$
in the respective lists (with rank $\infty$ if absent), and $k = 60$ is the
fusion constant.  Items not present in a list receive a contribution of zero
from that list.  The top-$K_r = 200$ items by RRF score form the candidate set
passed to the Prioritization Engine.

\paragraph{Cross-Encoder Reranking.}
A cross-encoder model jointly encodes each query-item pair to produce a
fine-grained relevance score that is incorporated into the priority function
(see \cref{subsec:prioritization}).  Cross-encoder inference is batched over
the $K_r$ candidates with a maximum batch size of 64.

\subsection{Prioritization Engine}
\label{subsec:prioritization}

The Prioritization Engine computes a scalar priority score $\prio(i, q)$ for
each candidate item $i$ given the current query $q$.  The priority function
integrates four dimensions:

\begin{equation}
  \prio(i, q) =
    w_r \cdot \rel(i, q)
  + w_e \cdot \rec(i)
  + w_m \cdot \alpha(i)
  + w_n \cdot \nov(i, \mathcal{S})
  \label{eq:priority}
\end{equation}

\noindent where the four components and their default weights are:

\paragraph{Relevance $\rel(i, q)$.}
The cross-encoder score $\text{CE}(q, i) \in [0, 1]$ (obtained from the
second-stage reranker) is used when available; otherwise, the normalized
cosine similarity $(\cos(\phi(i), \phi(q)) + 1) / 2$ is used.
Default weight: $w_r = 0.40$.

\paragraph{Recency $\rec(i)$.}
Temporal decay is modeled with an exponential decay function:
\begin{equation}
  \rec(i) = \exp\!\bigl(-\lambda \cdot \Delta t_i\bigr)
  \label{eq:recency}
\end{equation}
where $\Delta t_i = t_{\text{now}} - \tau(i)$ is the elapsed time in seconds
since item $i$ was created, and $\lambda = 10^{-4}\;\text{s}^{-1}$ corresponds
to a half-life of approximately 115 minutes—calibrated to the typical duration
of a long-horizon agent session.  Default weight: $w_e = 0.25$.

\paragraph{Importance $\alpha(i)$.}
The importance score $\alpha(i) \in [0,1]$ is assigned at item creation.
For items of type \textsc{goal} or \textsc{plan}, $\alpha(i) = 1.0$
unconditionally.  For tool results and observations, $\alpha(i)$ is set by
an automatic importance classifier that uses a lightweight 125M-parameter LM
to score the informational value of the content relative to the current goal.
Default weight: $w_m = 0.20$.

\paragraph{Novelty $\nov(i, \mathcal{S})$.}
To avoid filling the context with redundant information, the novelty score
penalizes items that are semantically similar to items already selected in
$\mathcal{S}$:
\begin{equation}
  \nov(i, \mathcal{S}) =
  \begin{cases}
    1 & \text{if } \mathcal{S} = \emptyset \\
    1 - \displaystyle\max_{j \in \mathcal{S}}
      \cos\!\bigl(\phi(i), \phi(j)\bigr) & \text{otherwise}
  \end{cases}
  \label{eq:novelty}
\end{equation}
This formulation follows the MMR principle \citep{carbonell1998mmr}: items that
are already well-represented in the selected set receive a lower novelty score
and are thus deprioritized.  Because $\nov(i, \mathcal{S})$ depends on the
current selection $\mathcal{S}$, the priority function is re-evaluated
incrementally as the scheduler adds items.  Default weight: $w_n = 0.15$.

The weight vector $(w_r, w_e, w_m, w_n) = (0.40, 0.25, 0.20, 0.15)$ was
selected via a grid search on the \contextbench{} validation set with grid
step $0.05$ subject to the constraint $\sum w_k = 1$.

\subsection{Context Scheduler}
\label{subsec:scheduler}

The Context Scheduler takes the priority-scored candidate set and selects a
budget-feasible subset.  It implements a greedy algorithm
(\cref{alg:greedy-scheduler}) that exploits the submodular structure of the
selection objective.

\begin{algorithm}[t]
\DontPrintSemicolon
\SetAlgoLined
\KwIn{Candidate set $\mathcal{I}$, query $q$, token budget $B$,
      priority function $\prio(\cdot, q)$}
\KwOut{Selected context set $\mathcal{S}$}
\BlankLine
$\mathcal{S} \leftarrow \emptyset$\;
$B_{\text{used}} \leftarrow 0$\;
Compute $\prio(i, q)$ for all $i \in \mathcal{I}$ using \cref{eq:priority}\;
Sort $\mathcal{I}$ in descending order of $\prio(i, q)$\;
\BlankLine
\ForEach{$i \in \mathcal{I}$ \emph{(in priority order)}}{
  \If{$B_{\text{used}} + \tokens(i) \leq B$}{
    $\mathcal{S} \leftarrow \mathcal{S} \cup \{i\}$\;
    $B_{\text{used}} \leftarrow B_{\text{used}} + \tokens(i)$\;
    Update $\nov(j, \mathcal{S})$ for all remaining $j \in \mathcal{I} \setminus \mathcal{S}$\;
    Re-sort remaining $\mathcal{I} \setminus \mathcal{S}$ by updated $\prio(j, q)$\;
  }
}
\BlankLine
Apply positional ordering: place \textsc{goal}/\textsc{plan} items first,
\textsc{observation}/\textsc{tool\_result} items last\;
\Return $\mathcal{S}$\;
\caption{\textsc{GreedyContextScheduler}}
\label{alg:greedy-scheduler}
\end{algorithm}

\paragraph{Approximation Guarantee.}
We now establish that the greedy scheduler achieves a provable approximation
ratio under a natural formalization of the objective.

\begin{theorem}[Greedy Approximation Ratio]
\label{thm:greedy}
Let $f(\mathcal{S}) = \sum_{i \in \mathcal{S}} \prio(i, q)$ be the priority
coverage objective (with the novelty component $\nov(i, \mathcal{S})$
evaluated incrementally as in \cref{alg:greedy-scheduler}).  When $f$ is
monotone submodular and $\tokens(i) = c$ for all $i$ (uniform token cost),
the greedy solution $\mathcal{S}_G$ satisfies:
\begin{equation}
  f(\mathcal{S}_G) \;\geq\; \Bigl(1 - \frac{1}{e}\Bigr) \cdot f(\mathcal{S}^*)
  \label{eq:greedy-bound}
\end{equation}
where $\mathcal{S}^* = \argmax_{|\mathcal{S}| \leq K} f(\mathcal{S})$ is the
optimal $K$-item solution and $K = \lfloor B / c \rfloor$.
\end{theorem}
\begin{proof}
The priority function $\prio(i, q)$ as defined in \cref{eq:priority} is
submodular in the selection set $\mathcal{S}$ because the novelty component
$\nov(i, \mathcal{S}) = 1 - \max_{j \in \mathcal{S}} \cos(\phi(i), \phi(j))$
is a coverage function: as $\mathcal{S}$ grows, $\nov(i, \mathcal{S})$ is
non-increasing for any fixed $i$ (adding more items to $\mathcal{S}$ can only
increase the maximum cosine similarity, weakly decreasing novelty).  The
sum over selected items of such a function is a non-negative linear combination
of submodular functions and is therefore submodular.  Under the uniform
token-cost assumption, the cardinality-constrained maximization of a monotone
submodular function by a greedy algorithm achieves the $(1 - 1/e)$ ratio
\citep{nemhauser1978analysis}.
\end{proof}

In the non-uniform-cost regime, we apply the Knapsack-greedy variant that
selects by the density $\prio(i, q) / \tokens(i)$, which is known to
achieve a $1/2$-approximation \citep{sviridenko2004note}.  In practice,
since \contextos{} operates on the priority-ranked list and breaks ties by
$\prio/\tokens$ density, empirical performance is substantially better than
either worst-case bound.

\subsection{Compression Engine}
\label{subsec:compression}

Even after greedy selection, the assembled context may exceed the available
token budget due to mandatory items (system prompt, current query, and items
of type \textsc{goal} that must always be included).  The Compression Engine
reduces item lengths while preserving semantic fidelity.

\paragraph{Compression Strategies.}
Three strategies are available, applied in order of increasing aggressiveness:
\begin{itemize}
  \item \textbf{Extractive compression:} sentences are scored by a sentence
  salience model (a 66M-parameter DistilBERT-based classifier) and the
  top-scoring sentences are retained.  Targets a compression ratio of $\rho \in
  [0.4, 0.8]$.
  \item \textbf{Abstractive compression:} a 1.5B-parameter summarization model
  (fine-tuned T5-XL) generates a condensed paraphrase of the item.  Targets
  $\rho \in [0.2, 0.5]$.  Used when the item is a long document or a multi-step
  tool response.
  \item \textbf{Hierarchical compression:} items are first segmented into
  passages, each passage is abstractively compressed, and the passage summaries
  are then concatenated.  Achieves $\rho \in [0.1, 0.35]$ while preserving more
  structural detail than flat abstractive summarization.
\end{itemize}

\paragraph{Progressive Compression.}
The strategy is selected based on the current budget pressure ratio
$\beta = B_{\text{used}} / B_{\max}$:
\begin{equation}
  \text{strategy}(\beta) =
  \begin{cases}
    \text{Extractive}   & \text{if } \beta \leq 0.85 \\
    \text{Abstractive}  & \text{if } 0.85 < \beta \leq 0.95 \\
    \text{Hierarchical} & \text{if } \beta > 0.95
  \end{cases}
  \label{eq:compression-strategy}
\end{equation}

\paragraph{Type-Conditional Compression.}
Items of type \textsc{goal} or \textsc{plan} are never compressed, as their
full content is essential for maintaining task coherence.  Items of type
\textsc{reasoning} are compressed only extractively (preserving logical
connectives and conclusions).  All other item types are eligible for any
compression strategy.

\paragraph{Quality Control.}
After compression, the semantic similarity between the original and compressed
representations is computed as $\text{sim}(i, \hat{i}) = \cos(\phi(i), \phi(\hat{i}))$.
If $\text{sim}(i, \hat{i}) < \theta_{\text{fidelity}} = 0.85$, the compression
is rejected and the original item is retained (potentially at the cost of
exceeding the token budget slightly; the scheduler accommodates up to 5\%
budget overflow in this case).

\subsection{Memory System}
\label{subsec:memory}

The \contextos{} Memory System implements a two-level hierarchy: Working Memory
and Long-Term Memory.

\paragraph{Working Memory.}
Working memory is a bounded priority heap with capacity $K_{\text{WM}} = 50$
items.  Items in working memory are those most recently accessed or most highly
prioritized.  The heap is keyed by priority $\prio(i, q_{\text{last}})$, where
$q_{\text{last}}$ is the most recent query.  Insertion into a full working
memory triggers the Governance Engine's eviction policy (see
\cref{subsec:governance}).

\paragraph{Long-Term Memory.}
Long-term memory is organized into three tiers:
\begin{itemize}
  \item \textbf{Episodic memory} stores event sequences: time-stamped
  observations, action records, and their outcomes.  Items in episodic memory
  are indexed by both timestamp and embedding for temporal range queries and
  semantic search.
  \item \textbf{Semantic memory} stores factual assertions extracted from
  observations and tool results.  Facts are represented as (subject, predicate,
  object) triples plus a confidence score, supporting structured queries in
  addition to semantic search.
  \item \textbf{Procedural memory} stores reusable action patterns—sequences of
  tool calls that successfully achieved a subgoal—indexed by the subgoal
  embedding.  At inference time, procedural memories provide executable plan
  templates that reduce the planning burden on the backbone LLM.
\end{itemize}

\paragraph{Promotion Criteria.}
An item in working memory is promoted to long-term memory when both of the
following conditions hold:
\begin{align}
  \text{access\_count}(i) &\geq 3 \label{eq:promotion-count} \\
  \alpha(i) &\geq 0.7 \label{eq:promotion-importance}
\end{align}
The combination of access frequency and importance ensures that only items that
are both repeatedly useful and inherently significant migrate to the more
persistent (and more expensive to query) long-term store.  The destination tier
within long-term memory is determined by $\text{type}(i)$: observations and
tool results go to episodic memory, extracted facts to semantic memory, and
successful action sequences to procedural memory.

\subsection{Governance Engine}
\label{subsec:governance}

The Governance Engine runs asynchronously in a background thread, executing
periodically (default: every 30 seconds or after every 100 agent steps,
whichever comes first) to enforce lifecycle policies across all memory tiers.

\paragraph{Retention Policies.}
Three retention policies are supported and can be composed:
\begin{itemize}
  \item \textbf{TimeBasedRetention($T$):} An item $i$ is retained if
  $\Delta t_i < T$.  Default: $T = 72\;\text{h}$ for working memory,
  $T = 720\;\text{h}$ for long-term episodic memory.
  \item \textbf{ImportanceBasedRetention($\theta$):} An item $i$ is retained
  if $\alpha(i) \geq \theta$.  Default: $\theta = 0.1$ (removes only very
  low-importance items).
  \item \textbf{FrequencyBasedRetention($f_{\min}$):} An item $i$ is retained
  if $\text{access\_count}(i) / \Delta t_i \geq f_{\min}$.  Default:
  $f_{\min} = 0.001\;\text{accesses/h}$.
\end{itemize}

An item is evicted when it fails \emph{any} active retention policy.  Before
eviction, the Governance Engine checks the promotion criteria in
\cref{eq:promotion-count,eq:promotion-importance}; items that meet both criteria
are promoted rather than evicted.

\paragraph{Forgetting.}
Importance scores are subject to exponential decay with a separate forgetting
rate $\lambda_f$:
\begin{equation}
  \alpha_t(i) = \alpha_0(i) \cdot \exp\!\bigl(-\lambda_f \cdot \Delta t_i\bigr)
  \label{eq:forgetting}
\end{equation}
with $\lambda_f = 5 \times 10^{-5}\;\text{s}^{-1}$ (half-life $\approx 3.9$ h).
This ensures that items that were once important but have not been accessed
recently gradually become candidates for eviction, mimicking the Ebbinghaus
forgetting curve \citep{ebbinghaus1885memory}.  Items of type \textsc{goal} are
exempt from forgetting ($\lambda_f = 0$ for \textsc{goal} items).

\paragraph{Compression Triggers.}
The Governance Engine also triggers preemptive compression before the scheduler
is forced into reactive compression.  A \textbf{TokenBudgetTrigger} fires when
the total token count of all working-memory items exceeds 80\% of $B_{\max}$,
initiating extractive compression on the lowest-priority quintile of working
memory items.

%% ============================================================
%% ============================================================
%%  SECTION 5: CONTEXTBENCH
%% ============================================================
%% ============================================================
\section{\contextbench{}: A Benchmark for Long-Horizon Context Management}
\label{sec:benchmark}

\subsection{Dataset Construction}
\label{subsec:dataset-construction}

\contextbench{} is a purpose-built evaluation benchmark for long-horizon context
management.  It consists of 70,000 tasks drawn from eight task types across six
domains, with context lengths ranging from 512 to 32,768 tokens and three
difficulty levels (easy, medium, hard).  The benchmark is partitioned into
60,000 training examples, 5,000 validation examples, and 5,000 held-out test
examples.  The test split is used for all reported results in this paper.

Tasks were constructed via a combination of expert curation and template-based
generation.  For each task, a \emph{ground-truth context} is defined as the
minimal subset of context items necessary and sufficient for a well-informed
agent to produce the correct answer.  Distractor items (semantically related
but irrelevant) are added to reach the target context length.  The difficulty
level is determined by the ratio of distractor items to ground-truth items
(easy: 2:1, medium: 5:1, hard: 10:1) and by the degree of semantic overlap
between distractor and ground-truth items.  Annotators (domain experts)
verified all ground-truth context labels and task answers via a two-round
agreement process; inter-annotator agreement (Cohen's $\kappa$) is $0.84$
across all task types.

The six domains are: (1) Scientific Literature, (2) Legal Reasoning,
(3) Software Engineering, (4) Financial Analysis, (5) Historical Research,
and (6) Medical Diagnosis.  Each domain contributes approximately equally to
the 5,000 test examples.

\subsection{Task Types}
\label{subsec:task-types}

\cref{tab:task-types} provides an overview of the eight task types in
\contextbench{}.

\begin{table}[t]
\centering
\caption{Task types in \contextbench{}.  ``Avg.\ Relevant Items'' refers to
  the mean number of ground-truth context items per task.
  ``Avg.\ Context Tokens'' is the mean total context size including distractors.}
\label{tab:task-types}
\begin{tabular}{llrrr}
\toprule
\textbf{Task Type} & \textbf{Description} & \textbf{Count} &
\textbf{Avg.\ Rel.\ Items} & \textbf{Avg.\ Ctx.\ Tokens} \\
\midrule
Multi-hop QA         & Chain of $\geq 3$ reasoning hops   & 11{,}500 & 5.2 & 8{,}342 \\
Causal Chain         & Root-cause and effect tracing       & 9{,}000  & 4.1 & 7{,}615 \\
Procedural Planning  & Multi-step action sequence planning & 9{,}500  & 6.8 & 11{,}200 \\
Code Debugging       & Error diagnosis across files        & 10{,}000 & 7.3 & 13{,}480 \\
Literature Synthesis & Cross-document claim synthesis      & 8{,}000  & 9.1 & 18{,}920 \\
Timeline Reconstruction & Temporal event ordering          & 7{,}500  & 5.9 & 9{,}870 \\
Cross-domain Analogy & Structural similarity detection     & 7{,}000  & 3.7 & 6{,}150 \\
Adversarial Distract.& Resistance to irrelevant context    & 7{,}500  & 2.8 & 14{,}300 \\
\midrule
\textbf{Total / Mean} & & \textbf{70{,}000} & \textbf{5.6} & \textbf{11{,}235} \\
\bottomrule
\end{tabular}
\end{table}

\subsection{Evaluation Metrics}
\label{subsec:metrics}

We evaluate all systems on four complementary metrics:

\paragraph{Task Success Rate (TSR).}
TSR is the primary metric.  For each task, the system's prediction is compared
to the ground-truth answer using token-level F1 (following the SQuAD evaluation
protocol \citep{rajpurkar2016squad}):
\begin{equation}
  \text{F1}_{\text{task}} = \frac{2 \cdot P \cdot R}{P + R}
  \label{eq:f1}
\end{equation}
where $P$ and $R$ are precision and recall over answer tokens after stopword
removal.  TSR is the mean $\text{F1}_{\text{task}}$ across all test examples.
A task is considered ``successful'' (for reporting purposes in qualitative
analyses) if $\text{F1}_{\text{task}} \geq 0.5$.

\paragraph{Context Relevance Score (CRS).}
CRS measures the quality of the context set selected by the system.  Given the
ground-truth relevant item set $\mathcal{G}$ and the system-selected set
$\mathcal{S}$, we report:
\begin{equation}
  \text{P@}K = \frac{|\mathcal{S}_{:K} \cap \mathcal{G}|}{K}, \qquad
  \text{NDCG@}10 = \frac{\text{DCG@}10}{\text{IDCG@}10}
  \label{eq:crs}
\end{equation}
where $\text{DCG@}10 = \sum_{k=1}^{10} \frac{\mathbbm{1}[i_k \in \mathcal{G}]}
{\log_2(k+1)}$ and IDCG@10 is the ideal DCG achieved when all top-10 positions
are filled with relevant items.

\paragraph{Token Efficiency (TE).}
Token Efficiency measures the mean number of tokens consumed per successfully
completed task (lower is better):
\begin{equation}
  \text{TE} = \frac{\sum_{t \in T_{\text{success}}} \tokens(\mathcal{C}_t)}
              {|T_{\text{success}}|}
  \label{eq:te}
\end{equation}
where $T_{\text{success}}$ is the set of tasks for which
$\text{F1}_{\text{task}} \geq 0.5$.

\paragraph{Latency.}
We report the wall-clock time for context preparation (retrieval through
context assembly, excluding LLM inference) in milliseconds, measured on a
single NVIDIA A100 80 GB GPU.

%% ============================================================
%% ============================================================
%%  SECTION 6: EXPERIMENTAL SETUP
%% ============================================================
%% ============================================================
\section{Experimental Setup}
\label{sec:setup}

\subsection{Baselines}
\label{subsec:baselines}

We compare \contextos{} against five baselines representing the principal
families of context management strategies:

\begin{enumerate}[label=\textbf{B\arabic*.}, leftmargin=3em]
  \item \textbf{Full Context (FC):}  The entire available context is passed to
  the LLM without any management.  When the total context exceeds the window
  limit, the task is marked as failed.  This baseline reveals the ceiling
  performance achievable with unlimited window capacity.

  \item \textbf{Truncation / First-$K$:}  The first $B_{\max}$ tokens of the
  context stream are retained and the remainder discarded.  This is the
  simplest and most widely deployed production strategy.

  \item \textbf{RAG-Only:}  A standard RAG pipeline using the same BGE-M3
  encoder and Qdrant backend as \contextos{}, but without the Prioritization
  Engine, Memory System, or Governance Engine.  Retrieves the top-$K$ items
  by cosine similarity and fills the context greedily.

  \item \textbf{MemGPT} \citep{packer2023memgpt}:  The virtual-context
  management system with in-context main memory and out-of-context archival
  memory, with function-call-based memory operations.

  \item \textbf{RAPTOR} \citep{sarthi2024raptor}:  Hierarchical summarization
  tree with multi-level retrieval.  Items are indexed in a tree of recursive
  abstractive summaries; retrieval traverses the tree to identify the most
  relevant leaf or internal nodes.

\end{enumerate}

\cref{tab:baselines-summary} summarizes the key properties of each system.

\begin{table}[t]
\centering
\caption{Summary of baseline systems and their context management capabilities.
  \checkmark = fully supported, $\circ$ = partially supported, $\times$ = not
  supported.}
\label{tab:baselines-summary}
\begin{tabular}{lccccc}
\toprule
\textbf{System} & \textbf{Retrieval} & \textbf{Compression} &
\textbf{Priority} & \textbf{Memory Hier.} & \textbf{Governance} \\
\midrule
Full Context  & $\times$    & $\times$    & $\times$    & $\times$    & $\times$ \\
Truncation    & $\times$    & $\times$    & $\times$    & $\times$    & $\times$ \\
RAG-Only      & \checkmark  & $\times$    & $\circ$     & $\times$    & $\times$ \\
MemGPT        & \checkmark  & $\circ$     & $\circ$     & $\circ$     & $\times$ \\
RAPTOR        & \checkmark  & \checkmark  & $\circ$     & $\circ$     & $\times$ \\
\textbf{\contextos{}} & \checkmark & \checkmark & \checkmark & \checkmark & \checkmark \\
\bottomrule
\end{tabular}
\end{table}

\subsection{Backbone Models}
\label{subsec:backbone}

All systems are evaluated with three backbone LLMs:
\begin{itemize}
  \item \textbf{GPT-OSS-20B:}  A 20-billion-parameter open-source model
  following the GPT architecture, pre-trained on a 2T-token multilingual
  corpus, with a 32K-token context window.
  \item \textbf{Qwen3:}  The third generation of the Qwen series
  \citep{bai2023qwen}, a strong multilingual model with 72B parameters
  (using 4-bit AWQ quantization for memory efficiency) and a 32K context window.
  \item \textbf{GLM-4.5:}  ChatGLM's fourth-generation model with a
  bilingual (Chinese-English) focus, 130B parameters (8-bit), and a 128K
  context window.
\end{itemize}

All backbone models are evaluated in a zero-shot setting (no fine-tuning or
few-shot examples specific to \contextbench{}) to isolate the effect of context
management rather than task-specific training.

\subsection{Implementation Details}
\label{subsec:implementation}

\textbf{Hardware:}  All experiments are run on a server with 8$\times$NVIDIA
A100 80 GB GPUs, 512 GB DDR4 RAM, and 2$\times$AMD EPYC 7763 CPUs.

\textbf{Embedding model:}  \texttt{BAAI/bge-m3} is the default dense encoder
(1024-d, 570M parameters).  We also evaluate \texttt{intfloat/e5-large-v2}
(1024-d, 335M parameters) in the retrieval ablation.

\textbf{Cross-encoder reranker:}  \texttt{cross-encoder/ms-marco-MiniLM-L-6-v2}
(22M parameters, latency 1.8 ms per query-item pair on A100).

\textbf{Vector store:}  Qdrant (in-process mode) with HNSW indices
($m = 16$, $\text{ef\_construction} = 200$).

\textbf{Compression models:}  Extractive—DistilBERT-based sentence scorer
(66M); Abstractive—T5-XL (3B) fine-tuned on a mixture of CNN/DailyMail and
XSum; Hierarchical—the same T5-XL applied per-passage.

\textbf{Token counting:}  All token budgets are measured with the backbone
model's tokenizer.  Default context budget: 8,192 tokens (for the 8K
experimental condition); experiments at 512, 4K, and 32K tokens use the
corresponding budget.

\textbf{Inference:}  LLM inference is performed with vLLM \citep{kwon2023vllm}
with continuous batching; batch size 16 for GPT-OSS-20B, 8 for Qwen3, 4 for
GLM-4.5.

\textbf{Reproducibility:}  All experiments use fixed random seeds
($\text{seed} = 42$).  Code, model checkpoints (for compression models),
and benchmark data will be released upon publication.

%% ============================================================
%% ============================================================
%%  SECTION 7: RESULTS AND ANALYSIS
%% ============================================================
%% ============================================================
\section{Results and Analysis}
\label{sec:results}

\subsection{Main Results}
\label{subsec:main-results}

\cref{tab:main-results} reports the Task Success Rate (TSR, mean $\pm$ std
over three random seeds) for all six methods across four context-length regimes
on the \contextbench{} test set, averaged over the three backbone LLMs.

\begin{table}[t]
\centering
\caption{Task Success Rate (\%, mean $\pm$ std) on \contextbench{} test set,
  averaged over GPT-OSS-20B, Qwen3, and GLM-4.5.  Bold indicates the best
  result per context-length column.  All \contextos{} improvements over MemGPT
  and RAPTOR are statistically significant ($p < 0.001$, paired $t$-test with
  Bonferroni correction).}
\label{tab:main-results}
\resizebox{\linewidth}{!}{%
\begin{tabular}{lcccc}
\toprule
\multirow{2}{*}{\textbf{Method}} &
\multicolumn{4}{c}{\textbf{Context Length (tokens)}} \\
\cmidrule(lr){2-5}
 & \textbf{512} & \textbf{4K} & \textbf{8K} & \textbf{32K} \\
\midrule
Full Context     & $81.1 \pm 0.8$ & $79.3 \pm 0.9$ & $74.6 \pm 1.1$ & $68.2 \pm 1.3$ \\
Truncation       & $78.9 \pm 0.7$ & $61.4 \pm 1.2$ & $41.8 \pm 1.5$ & $31.5 \pm 1.6$ \\
RAG-Only         & $79.4 \pm 0.9$ & $74.2 \pm 1.0$ & $65.1 \pm 1.3$ & $48.7 \pm 1.4$ \\
MemGPT           & $80.2 \pm 0.8$ & $77.1 \pm 1.0$ & $70.4 \pm 1.2$ & $58.9 \pm 1.4$ \\
RAPTOR           & $79.8 \pm 0.9$ & $76.3 \pm 1.1$ & $68.9 \pm 1.3$ & $54.1 \pm 1.5$ \\
\textbf{\contextos{}} & $\mathbf{84.2 \pm 0.6}$ & $\mathbf{82.9 \pm 0.7}$ & $\mathbf{79.6 \pm 0.9}$ & $\mathbf{71.3 \pm 1.1}$ \\
\midrule
$\Delta$ vs.\ MemGPT      & $+4.0$ & $+5.8$ & $+9.2$ & $+12.4$ \\
$\Delta$ vs.\ Truncation  & $+5.3$ & $+21.5$ & $+37.8$ & $+39.8$ \\
\bottomrule
\end{tabular}%
}
\end{table}

\paragraph{ContextOS achieves the highest TSR at all context lengths.}
At 512 tokens, the performance gap between methods is modest (all methods
achieve 78.9–84.2\%), as the context is short enough that simple strategies
suffice for easy tasks.  However, as context length grows to 4K, 8K, and
32K tokens, the performance differential widens dramatically.  \contextos{}
achieves 71.3\% at 32K tokens—a 12.4 percentage-point (pp) improvement over
MemGPT ($p < 0.001$, Cohen's $d = 0.91$), a 17.2 pp improvement over RAPTOR,
and a 39.8 pp improvement over na\"{i}ve truncation.

\paragraph{Truncation degrades most severely with scale.}
Truncation falls from 78.9\% at 512 tokens to only 31.5\% at 32K tokens—a
decline of 47.4 pp.  This collapse reflects the fundamental inadequacy of
recency-only retention: as context length grows, an increasing fraction of the
relevant items are old and therefore discarded.  The Adversarial Distraction and
Timeline Reconstruction task types are particularly impacted, as they require
integrating information spread across the full context length.

\paragraph{\contextos{} degrades most gracefully.}
Across the four context-length conditions, \contextos{}'s TSR decreases by only
12.9 pp (from 84.2\% at 512 to 71.3\% at 32K), compared to 22.3 pp for MemGPT
and 47.4 pp for Truncation.  This graceful degradation is attributable to the
combined effect of the priority scheduler (which selects the most relevant
subset even when the context is highly overloaded) and the long-term memory
governance (which retains historically important items even after they have
aged out of working memory).

\paragraph{Token efficiency.}
\contextos{} uses a mean of 3,834 tokens per successful task at the 8K
condition, compared to 10,315 tokens for Full Context—a 62.9\% reduction.
This efficiency gain is particularly significant for API-based deployment
scenarios, where token consumption directly drives inference cost.  RAG-Only
achieves similar token efficiency (3,902 tokens) but at substantially lower
TSR, confirming that token reduction alone is not sufficient; the quality of
context selection matters.

\subsection{Retrieval Quality}
\label{subsec:retrieval-quality}

\cref{tab:retrieval-quality} reports retrieval quality metrics for the
retrieval-capable baselines on the \contextbench{} test set.

\begin{table}[t]
\centering
\caption{Retrieval quality on \contextbench{} test set.  P@1, P@5, and
  NDCG@10 measure how well each system's selected context matches the
  ground-truth relevant items.  Higher is better.}
\label{tab:retrieval-quality}
\begin{tabular}{lccc}
\toprule
\textbf{System} & \textbf{P@1} & \textbf{P@5} & \textbf{NDCG@10} \\
\midrule
RAG-Only (Dense)          & 0.614 & 0.571 & 0.623 \\
RAG-Only (Sparse)         & 0.582 & 0.541 & 0.597 \\
RAG-Only (Hybrid, no rerank) & 0.651 & 0.608 & 0.657 \\
MemGPT                    & 0.672 & 0.631 & 0.683 \\
RAPTOR                    & 0.658 & 0.619 & 0.671 \\
\textbf{\contextos{}}     & \textbf{0.741} & \textbf{0.703} & \textbf{0.762} \\
\bottomrule
\end{tabular}
\end{table}

\contextos{} achieves NDCG@10 = 0.762, substantially outperforming all
baselines.  The improvement over RAG-Only with hybrid retrieval but no
reranking (NDCG@10 = 0.657) demonstrates the value of the cross-encoder
second stage, which increases NDCG@10 by 0.105 absolute.  The further gain
over MemGPT (NDCG@10 = 0.683) reflects the contribution of the priority
scheduler's MMR-based novelty component, which improves the diversity of
the selected context set and reduces redundancy.  Sparse retrieval alone
underperforms dense retrieval by 2.6 pp in NDCG@10, but the hybrid
combination improves over dense alone by 3.4 pp, confirming the complementary
nature of dense and sparse signals.

\subsection{Compression Quality}
\label{subsec:compression-quality}

We evaluate the three compression strategies on a random 500-item sample from
\contextbench{} test tasks, measuring compression ratio $\rho$,
semantic fidelity (cosine similarity between original and compressed embeddings),
and downstream TSR impact.  Results are shown in \cref{fig:compression-quality}
(referenced here; see the submitted supplementary materials).  Hierarchical
compression achieves the best trade-off: at a compression ratio of $\rho = 0.4$,
it maintains a semantic fidelity of 0.91 and a TSR within 1.3 pp of the
uncompressed baseline.  Abstractive compression achieves higher compression
($\rho = 0.31$) but at a fidelity cost of 0.87 and a TSR penalty of 2.8 pp.
Extractive compression is most conservative ($\rho = 0.62$) but preserves
fidelity almost perfectly (0.97) and incurs only a 0.4 pp TSR penalty.

These results validate the progressive compression strategy of \cref{eq:compression-strategy}:
extractive compression is preferable when budget pressure is mild (preserving
fidelity), while hierarchical compression is appropriate under severe budget
pressure (achieving maximum space savings at acceptable fidelity cost).
We recommend a default compression ratio of $\rho = 0.4$ using hierarchical
compression for \contextos{} deployments, as it provides the best balance of
token savings and task performance across the \contextbench{} task types.

\subsection{Latency Analysis}
\label{subsec:latency}

\cref{tab:latency} reports the wall-clock latency breakdown of the \contextos{}
pipeline on the A100 hardware.

\begin{table}[t]
\centering
\caption{Latency breakdown of \contextos{} pipeline components (milliseconds,
  mean $\pm$ std over 1,000 test tasks, excluding LLM inference time).}
\label{tab:latency}
\begin{tabular}{lrr}
\toprule
\textbf{Component} & \textbf{Mean (ms)} & \textbf{Std (ms)} \\
\midrule
Dense Retrieval (BGE-M3 + HNSW)  & 18.3 & 2.1 \\
Sparse Retrieval (BM25)           & 6.4  & 0.8 \\
RRF Fusion                        & 2.1  & 0.2 \\
Cross-Encoder Reranking           & 15.3 & 3.4 \\
\textit{Retrieval subtotal}       & \textit{42.1} & \textit{4.8} \\
Priority Scoring                  & 12.7 & 1.6 \\
Greedy Scheduling                 & 4.3  & 0.5 \\
\textit{Scheduling subtotal}      & \textit{17.0} & \textit{1.9} \\
Compression (when triggered)      & 31.4 & 8.2 \\
Memory Registry Lookup            & 8.6  & 1.0 \\
Context Assembly                  & 7.1  & 0.9 \\
\textit{Total (no compression)}   & \textit{124.8} & \textit{14.2} \\
\textit{Total (with compression)} & \textit{156.2} & \textit{22.4} \\
\bottomrule
\end{tabular}
\end{table}

The total pipeline latency of 156.2 ms (when compression is triggered) is well
within interactive response time tolerances (typically 200--500 ms for
conversational applications).  Retrieval dominates the latency budget at
42.1 ms (27.0\% of total), with the cross-encoder reranking step being the
primary bottleneck within retrieval.  Compression, when triggered (at 8K and
32K conditions), contributes 31.4 ms.  The scheduling components (priority
scoring and greedy selection) are inexpensive at 17.0 ms combined.  These
results confirm that \contextos{}'s performance gains come without prohibitive
latency overhead.

%% ============================================================
%% ============================================================
%%  SECTION 8: ABLATION STUDY
%% ============================================================
%% ============================================================
\section{Ablation Study}
\label{sec:ablation}

To assess the individual contribution of each \contextos{} component, we
conduct a systematic ablation study at the 8K context length condition
(the primary operating regime), using GPT-OSS-20B as the backbone model.
We evaluate seven variants, each corresponding to removing or disabling one
component, as well as a full system variant.  Results are reported in
\cref{tab:ablation}.

\begin{table}[t]
\centering
\caption{Ablation study results.  TSR (\%) on \contextbench{} test set
  at 8K context length, GPT-OSS-20B backbone.  $\Delta$\,TSR is the difference
  from the full \contextos{} system (negative = worse).  ``LT Memory'' = Long-Term
  Memory tier disabled (working memory only).  ``No Govern.'' = Governance
  Engine disabled.  ``No Compress.'' = Compression Engine disabled.
  ``Flat Sched.'' = Priority Scheduler replaced by uniform random sampling.
  ``No MMR'' = Novelty component ($w_n = 0$, weight redistributed to $w_r$).
  ``No Decay'' = Recency decay disabled ($\lambda = 0$).
  ``Dense Only'' = Sparse retrieval disabled.}
\label{tab:ablation}
\begin{tabular}{lcc}
\toprule
\textbf{System Variant} & \textbf{TSR (\%)} & \textbf{$\Delta$ TSR (pp)} \\
\midrule
\textbf{\contextos{} (Full)}  & \textbf{79.6} & --- \\
\midrule
w/o Long-Term Memory (LT)     & 63.2 & $-16.4$ \\
w/o Governance Engine         & 64.6 & $-15.0$ \\
w/o Compression Engine        & 73.0 & $-6.6$  \\
Flat Scheduler (random)       & 66.8 & $-12.8$ \\
w/o MMR Novelty (No MMR)      & 75.1 & $-4.5$  \\
w/o Recency Decay (No Decay)  & 76.9 & $-2.7$  \\
Dense Retrieval Only          & 77.3 & $-2.3$  \\
\midrule
Sum of individual ablations   &  ---  & $-60.3$ \\
Actual total gain vs.\ RAG-Only & --- & $+14.5$ \\
Interactions (synergy)        &  ---  & $-22.7$ \\
\bottomrule
\end{tabular}
\end{table}

\paragraph{Long-Term Memory is the most critical component.}
Disabling the long-term memory tier (forcing all items into working memory only,
with capacity $K_{\text{WM}} = 50$) yields the largest single-component drop:
$-16.4$ pp TSR.  Without long-term memory, items that have not been recently
accessed but remain relevant to the task (e.g., early goal specifications,
domain facts retrieved in previous agent steps) are evicted from working
memory and are no longer available for scheduling.  This finding validates
the central thesis of \contextos{}: that a hierarchical memory architecture
with explicit lifecycle management is essential for long-horizon performance.

\paragraph{Governance Engine is the second most critical component.}
Removing governance ($-15.0$ pp) prevents the memory system from promoting
high-value items from working to long-term memory and from enforcing
retention policies.  Without governance, the working memory fills with
low-importance, recently-created items while high-importance older items are
lost.  The similarity of the governance and long-term memory ablation results
reflects the tight coupling between these two components: the long-term memory
tier only provides value if the governance engine correctly populates it.

\paragraph{Priority Scheduler accounts for $-12.8$ pp.}
Replacing the priority scheduler with random selection from the retrieved
candidate set ($-12.8$ pp) demonstrates that the specific multi-factor
priority function of \cref{eq:priority} is substantially more effective
than unstructured selection.  This ablation controls for retrieval quality
(the same hybrid retrieval is used), isolating the contribution of the
scheduling policy itself.

\paragraph{Compression Engine contributes $-6.6$ pp.}
The contribution of the Compression Engine is most apparent at high context
lengths: at 32K tokens, removing compression results in a $-9.2$ pp drop
(not shown in the table), compared to $-6.6$ pp at 8K.  At 512 tokens,
the compression engine is rarely triggered, and its ablation has negligible
effect ($-0.3$ pp).  This confirms that compression is most valuable as
a mechanism for graceful degradation under severe budget pressure.

\paragraph{MMR Novelty and Recency Decay provide moderate gains.}
Removing the MMR novelty component ($w_n = 0$, $-4.5$ pp) and the recency
decay ($\lambda = 0$, $-2.7$ pp) both reduce performance, but less severely
than the structural components.  The MMR ablation is particularly impactful
on Literature Synthesis tasks (where context items are often semantically
similar across multiple retrieved documents), while the recency decay ablation
most strongly affects Timeline Reconstruction tasks (where temporal ordering
is critical).

\paragraph{Hybrid retrieval provides a modest but consistent gain.}
Disabling sparse retrieval ($-2.3$ pp) confirms the value of the hybrid
approach, consistent with the retrieval quality results in
\cref{subsec:retrieval-quality}.  The gain is consistent across all task types,
with the largest benefit on Adversarial Distraction tasks (where lexical
keyword matching is effective at identifying distractor items for exclusion).

\paragraph{Component Interactions.}
The sum of individual ablation penalties (60.3 pp) exceeds the actual total
gain of \contextos{} over RAG-Only (14.5 pp, at 8K), with the difference
(22.7 pp, or 37.6\% of the summed individual contributions) attributable to
synergistic interactions between components.  In particular, the long-term
memory tier and the governance engine are strongly synergistic: neither
provides full value without the other, as the governance engine's promotion
policies exist specifically to populate the long-term tier with high-value
items.  Similarly, compression and scheduling interact positively: a tighter
token budget (after compression) is most effectively utilized by the priority
scheduler, which allocates budget to the highest-value items rather than
filling it uniformly.

%% ============================================================
%% ============================================================
%%  SECTION 9: DISCUSSION
%% ============================================================
%% ============================================================
\section{Discussion}
\label{sec:discussion}

\subsection{When Does \contextos{} Help Most?}
\label{subsec:when-helps}

The results in \cref{sec:results,sec:ablation} make clear that \contextos{}
provides the greatest benefit in two regimes.  First, at long context lengths
(8K tokens and above), where the discrepancy between available context capacity
and the amount of information generated during agent execution is large.  In
this regime, the priority scheduler, long-term memory, and governance engine
work together to select and preserve the most relevant items, preventing the
catastrophic information loss that afflicts truncation and the quality
degradation that affects simpler retrieval-based methods.  The 39.8 pp gain
over truncation at 32K tokens is the clearest illustration of this benefit.

Second, \contextos{} provides disproportionate gains on multi-session tasks
and tasks with long causal chains.  In these settings, relevant information
from early in the task history must remain accessible many steps later.
The long-term memory hierarchy—particularly episodic memory for event sequences
and semantic memory for extracted facts—provides a structured repository for
this information that is indexed for rapid retrieval when needed.  Single-tier
working memory systems, including MemGPT in the evaluated configuration, are
insufficient for tasks that require more than $K_{\text{WM}} = 50$ items to be
simultaneously accessible.

In contrast, \contextos{}'s advantages are modest at short context lengths
(512 tokens) and for simple tasks (easy difficulty in \contextbench{}).  At
these scales, most baselines perform adequately, and the overhead of the
\contextos{} pipeline (primarily retrieval and cross-encoder latency) is not
justified relative to simpler alternatives.  Practitioners should consider
deploying lightweight baselines (e.g., RAG-Only) for short-context or
low-complexity settings.

\subsection{Limitations}
\label{subsec:limitations}

\textbf{Latency overhead.}  The \contextos{} pipeline adds 124.8--156.2 ms
of pre-LLM processing latency.  While this is acceptable for most applications,
it may be prohibitive for ultra-low-latency settings (e.g., sub-50 ms
voice assistants).  Future work should explore asynchronous pre-fetching and
speculative context preparation to reduce effective latency.

\textbf{Dependency on embedding quality.}  The retrieval, prioritization, and
novelty components all rely on dense embeddings from BGE-M3.  In domains where
BGE-M3's pre-training distribution is poorly matched (e.g., highly specialized
technical domains, low-resource languages), retrieval quality may degrade,
cascading through the pipeline.  Domain-adaptive fine-tuning of the encoder
or ensemble embedding strategies may mitigate this.

\textbf{Governance policy tuning.}  The retention thresholds, forgetting rates,
and promotion criteria in the Governance Engine were tuned on the \contextbench{}
validation set.  These hyperparameters may require re-tuning for deployment
in domains not represented in \contextbench{}, particularly for tasks with
very different temporal dynamics (e.g., batch processing vs.\ interactive
sessions) or importance distributions.

\textbf{Fixed memory capacity.}  The working memory capacity
$K_{\text{WM}} = 50$ is a fixed parameter in the current implementation.
Future work should explore dynamic capacity adjustment based on task complexity
and available computational budget.

\textbf{Evaluation scope.}  \contextbench{} covers eight task types and six
domains, but cannot comprehensively cover all long-horizon agent applications.
In particular, physical world interaction (robotics, embodied AI), real-time
streaming contexts, and adversarial environments are not included.  Expanding
\contextbench{} to these settings is an important direction for future work.

\subsection{Broader Implications}
\label{subsec:implications}

\contextos{} suggests that \emph{context engineering}—the systematic design
of policies for what information enters, persists in, and exits an agent's
context window—deserves recognition as a first-class research problem
within the AI systems community.  Just as the development of virtual memory
and page replacement algorithms was essential for scaling computer applications
beyond the limits of physical RAM, principled context lifecycle management
will likely be essential for scaling autonomous agents beyond the limits of
current context windows.

The OS abstraction is not merely metaphorical.  The formal parallels are
tight: context windows correspond to physical RAM, long-term memory tiers
correspond to storage hierarchy levels, the priority scheduler corresponds
to an OS process scheduler, and the governance engine corresponds to a
combination of a garbage collector and a cache replacement policy.  This
correspondence suggests that decades of OS and systems research may be directly
applicable to agent context management, and that the reverse may also be true:
the richer semantic structure of language-based context items (compared to
memory pages) may inspire new policies for classical memory management.

At a practical level, \contextos{} provides a deployable platform for
long-horizon agent applications.  Its modular architecture allows individual
components to be replaced or upgraded independently—for example, substituting
a stronger retrieval model or a more efficient compression model—without
requiring changes to the overall lifecycle management logic.

\subsection{Ethical Considerations}
\label{subsec:ethics}

The memory governance mechanisms in \contextos{} have direct implications
for user privacy.  An agent with a long-term memory system that persists
information across sessions can accumulate sensitive personal data (medical
information, financial details, private communications) that may pose privacy
risks if the memory is not properly governed.  The Governance Engine's
retention and forgetting policies provide a technical foundation for privacy
compliance: time-based retention policies can implement data minimization
principles, importance-based retention can prioritize task-relevant information
over incidentally sensitive details, and explicit eviction can implement
``right to be forgotten'' requests.  We recommend that deployments of \contextos{}
implement user-controlled memory governance dashboards that give users visibility
into and control over what the agent remembers.  Future work should also
investigate differential privacy mechanisms for the vector store and the
importance classifier to prevent membership inference attacks against agent
memories.

%% ============================================================
%% ============================================================
%%  SECTION 10: CONCLUSION
%% ============================================================
%% ============================================================
\section{Conclusion}
\label{sec:conclusion}

We have presented \contextos{}, an operating-system-inspired framework for
context lifecycle management in long-horizon autonomous agents.  \contextos{}
treats the context window as a managed resource, applying scheduling,
caching, compression, and eviction policies across a six-component pipeline:
a Retrieval Engine with hybrid dense-sparse retrieval and cross-encoder
reranking, a Prioritization Engine with a multi-factor priority function
incorporating temporal decay and MMR-based novelty, a Context Scheduler with
a provable submodularity guarantee, a progressive Compression Engine with
type-conditional policies, a two-level Memory System with episodic, semantic,
and procedural long-term tiers, and a Governance Engine enforcing retention,
forgetting, and promotion.

On \contextbench{}, a new benchmark of 70,000 long-horizon agent tasks,
\contextos{} achieves a task-success rate of 71.3\% at 32K tokens—a
12.4 pp improvement over MemGPT, 17.2 pp over RAPTOR, and 39.8 pp over
na\"{i}ve truncation—while reducing token consumption by 62.9\% relative to
Full Context.  The system operates with a total context-preparation latency
of 156.2 ms, well within interactive application tolerances.  Ablation studies
confirm that every component contributes positively, with long-term memory
governance and the priority scheduler contributing the largest individual gains.

Future work will focus on three directions.  First, \textit{adaptive policy
learning}: using reinforcement learning or Bayesian optimization to adapt
retention thresholds, decay rates, and compression ratios online during agent
deployment, rather than fixing them via offline hyperparameter search.
Second, \textit{cross-agent memory sharing}: extending the memory architecture
to multi-agent settings where multiple agents share a common semantic memory
while maintaining private episodic memories.  Third, \textit{privacy-preserving
memory}: integrating differential privacy mechanisms into the vector store and
the importance classifier to provide formal privacy guarantees for sensitive
user contexts.  We believe that context lifecycle management, as operationalized
by \contextos{}, represents a fundamental building block for the next generation
of reliable, efficient, and trustworthy long-horizon AI agents.

%% ============================================================
%%  ACKNOWLEDGEMENTS  (blinded for review)
%% ============================================================
\section*{Acknowledgements}
Omitted for double-blind review.

%% ============================================================
%%  DECLARATIONS
%% ============================================================
\section*{Data Availability}
The \contextbench{} dataset, \contextos{} source code, pre-trained compression
models, and evaluation scripts will be made publicly available at
\texttt{https://anonymous-url.github.io/contextos} upon acceptance.

\section*{Competing Interests}
The authors declare no competing interests.

%% ============================================================
%%  REFERENCES
%% ============================================================
\bibliographystyle{elsarticle-num-names}
\bibliography{references}

%% ============================================================
%%  INLINE REFERENCE LIST (for self-contained review submission)
%% ============================================================
\begin{thebibliography}{99}

\bibitem[Anderson and Lebiere(2004)]{anderson2004actr}
Anderson, J.~R., and Lebiere, C. (2004).
\newblock The atomic components of thought.
\newblock \textit{Lawrence Erlbaum Associates}.

\bibitem[AutoGPT(2023)]{autogpt2023}
AutoGPT (2023).
\newblock Auto-GPT: An autonomous GPT-4 experiment.
\newblock \texttt{https://github.com/Significant-Gravitas/AutoGPT}.

\bibitem[Bai et~al.(2023)]{bai2023qwen}
Bai, J., Bai, S., Chu, Y., Cui, Z., Dang, K., Deng, X., et~al. (2023).
\newblock Qwen technical report.
\newblock \textit{arXiv preprint arXiv:2309.16609}.

\bibitem[Beltagy et~al.(2020)]{beltagy2020longformer}
Beltagy, I., Peters, M.~E., and Cohan, A. (2020).
\newblock Longformer: The long-document transformer.
\newblock \textit{arXiv preprint arXiv:2004.05150}.

\bibitem[Carbonell and Goldstein(1998)]{carbonell1998mmr}
Carbonell, J., and Goldstein, J. (1998).
\newblock The use of MMR, diversity-based reranking for reordering documents
  and producing summaries.
\newblock In \textit{Proceedings of SIGIR}, pages 335--336.

\bibitem[Chen et~al.(2022)]{chen2022hybrid}
Chen, J., Gao, L., Hui, B., Lin, J., Gao, L., et~al. (2022).
\newblock Out-of-domain semantics to the rescue: Zero-shot hybrid retrieval
  models.
\newblock In \textit{Proceedings of ECIR}.

\bibitem[Chen et~al.(2024)]{chen2024bgem3}
Chen, J., Xiao, S., Zhang, P., Luo, K., Lian, D., and Liu, Z. (2024).
\newblock BGE M3-embedding: Multi-lingual, multi-functionality, multi-granularity
  text embeddings through self-knowledge distillation.
\newblock \textit{arXiv preprint arXiv:2402.03216}.

\bibitem[Chevalier et~al.(2023)]{chevalier2023autocompressors}
Chevalier, A., Wettig, A., Ajith, A., and Chen, D. (2023).
\newblock Adapting language models to compress contexts.
\newblock In \textit{Proceedings of EMNLP}, pages 3829--3846.

\bibitem[Cormack et~al.(2009)]{cormack2009rrf}
Cormack, G.~V., Clarke, C. L.~A., and Buettcher, S. (2009).
\newblock Reciprocal rank fusion outperforms Condorcet and individual rank
  learning methods.
\newblock In \textit{Proceedings of SIGIR}, pages 758--759.

\bibitem[Dao(2023)]{dao2023flashattention2}
Dao, T. (2023).
\newblock FlashAttention-2: Faster attention with better parallelism and work
  partitioning.
\newblock In \textit{Proceedings of ICLR 2024}.

\bibitem[Dao et~al.(2022)]{dao2022flashattention}
Dao, T., Fu, D., Ermon, S., Rudra, A., and Ré, C. (2022).
\newblock FlashAttention: Fast and memory-efficient exact attention with
  IO-awareness.
\newblock In \textit{Proceedings of NeurIPS}, pages 16344--16359.

\bibitem[Ebbinghaus(1885)]{ebbinghaus1885memory}
Ebbinghaus, H. (1885).
\newblock \textit{{\"U}ber das Ged{\"a}chtnis: Untersuchungen zur
  experimentellen Psychologie}.
\newblock Duncker \& Humblot, Leipzig.

\bibitem[Formal et~al.(2021)]{formal2021splade}
Formal, T., Piwowarski, B., and Clinchant, S. (2021).
\newblock SPLADE: Sparse lexical and expansion model for first stage ranking.
\newblock In \textit{Proceedings of SIGIR}, pages 2288--2292.

\bibitem[Garey and Johnson(1979)]{garey1979computers}
Garey, M.~R., and Johnson, D.~S. (1979).
\newblock \textit{Computers and Intractability: A Guide to the Theory of
  NP-Completeness}.
\newblock W. H. Freeman, New York.

\bibitem[Ge et~al.(2023)]{ge2023context}
Ge, T., Hu, J., Wang, L., Wang, B., Chen, S., Liu, T., and Wei, F. (2023).
\newblock In-context autoencoder for context compression in a large language
  model.
\newblock In \textit{Proceedings of ICLR 2024}.

\bibitem[Guu et~al.(2020)]{guu2020realm}
Guu, K., Lee, K., Tung, Z., Pasupat, P., and Chang, M.~W. (2020).
\newblock REALM: Retrieval-augmented language model pre-training.
\newblock In \textit{Proceedings of ICML}, pages 3929--3938.

\bibitem[Jiang et~al.(2023)]{jiang2023llmlingua}
Jiang, H., Wu, Q., Lin, C.-Y., Yang, Y., and Luo, L. (2023).
\newblock LLMLingua: Compressing prompts for accelerated inference of large
  language models.
\newblock In \textit{Proceedings of EMNLP}, pages 13358--13376.

\bibitem[Karpukhin et~al.(2020)]{karpukhin2020dpr}
Karpukhin, V., O{\v{g}}uz, B., Min, S., Lewis, P., Wu, L., Edunov, S.,
  Chen, D., and Yih, W.-T. (2020).
\newblock Dense passage retrieval for open-domain question answering.
\newblock In \textit{Proceedings of EMNLP}, pages 6769--6781.

\bibitem[Khattab and Zaharia(2020)]{khattab2020colbert}
Khattab, O., and Zaharia, M. (2020).
\newblock ColBERT: Efficient and effective passage search via contextualized
  late interaction over BERT.
\newblock In \textit{Proceedings of SIGIR}, pages 39--48.

\bibitem[Kim et~al.(2022)]{kim2022learned}
Kim, S., Shim, H., and Kim, Y. (2022).
\newblock Learned token pruning for transformers.
\newblock In \textit{Proceedings of KDD}, pages 784--794.

\bibitem[Kwon et~al.(2023)]{kwon2023vllm}
Kwon, W., Li, Z., Zhuang, S., Sheng, Y., Zheng, L., Yu, C.~H., Gonzalez,
  J.~E., Zhang, H., and Stoica, I. (2023).
\newblock Efficient memory management for large language model serving with
  PagedAttention.
\newblock In \textit{Proceedings of SOSP}, pages 611--626.

\bibitem[Laird(2012)]{laird2012soar}
Laird, J.~E. (2012).
\newblock \textit{The Soar Cognitive Architecture}.
\newblock MIT Press, Cambridge, MA.

\bibitem[Lewis et~al.(2020)]{lewis2020rag}
Lewis, P., Perez, E., Piktus, A., Petroni, F., Karpukhin, V., Goyal, N.,
  K{\"u}ttler, H., Lewis, M., Yih, W.-T., Rockt{\"a}schel, T., et~al. (2020).
\newblock Retrieval-augmented generation for knowledge-intensive NLP tasks.
\newblock In \textit{Proceedings of NeurIPS}, pages 9459--9474.

\bibitem[Liu et~al.(2023)]{liu2023lost}
Liu, N.~F., Lin, K., Hewitt, J., Paranjape, A., Bevilacqua, M., Petroni, F.,
  and Liang, P. (2023).
\newblock Lost in the middle: How language models use long contexts.
\newblock \textit{Transactions of the Association for Computational
  Linguistics}, 12:157--173.

\bibitem[Ma et~al.(2022)]{ma2022hybrid}
Ma, X., Wang, L., Yang, N., Wei, F., and Lin, J. (2022).
\newblock Fine-tuning LLaMA for multi-stage text retrieval.
\newblock \textit{arXiv preprint arXiv:2310.08319}.

\bibitem[Nemhauser et~al.(1978)]{nemhauser1978analysis}
Nemhauser, G.~L., Wolsey, L.~A., and Fisher, M.~L. (1978).
\newblock An analysis of approximations for maximizing submodular set functions.
\newblock \textit{Mathematical Programming}, 14(1):265--294.

\bibitem[Nogueira and Cho(2019)]{nogueira2019passage}
Nogueira, R., and Cho, K. (2019).
\newblock Passage re-ranking with BERT.
\newblock \textit{arXiv preprint arXiv:1901.04085}.

\bibitem[Packer et~al.(2023)]{packer2023memgpt}
Packer, C., Fang, V., Patil, S.~G., Lin, K., Wooders, S., and Gonzalez, J.~E.
  (2023).
\newblock MemGPT: Towards LLMs as operating systems.
\newblock \textit{arXiv preprint arXiv:2310.08560}.

\bibitem[Park et~al.(2023)]{park2023generative}
Park, J.~S., O'Brien, J.~C., Cai, C.~J., Morris, M.~R., Liang, P., and
  Bernstein, M.~S. (2023).
\newblock Generative agents: Interactive simulacra of human behavior.
\newblock In \textit{Proceedings of UIST}, pages 1--22.

\bibitem[Parvez et~al.(2021)]{parvez2021retrieval}
Parvez, M.~R., Ahmad, W.~U., Chakraborty, S., Ray, B., and Chang, K.-W.
  (2021).
\newblock Retrieval augmented code generation and summarization.
\newblock In \textit{Findings of EMNLP}, pages 2719--2734.

\bibitem[Press et~al.(2022)]{press2022alibi}
Press, O., Smith, N.~A., and Lewis, M. (2022).
\newblock Train short, test long: Attention with linear biases enables input
  length extrapolation.
\newblock In \textit{Proceedings of ICLR}.

\bibitem[Rajpurkar et~al.(2016)]{rajpurkar2016squad}
Rajpurkar, P., Zhang, J., Lopyrev, K., and Liang, P. (2016).
\newblock SQuAD: 100,000+ questions for machine comprehension of text.
\newblock In \textit{Proceedings of EMNLP}, pages 2383--2392.

\bibitem[Robertson and Zaragoza(2009)]{robertson2009bm25}
Robertson, S., and Zaragoza, H. (2009).
\newblock The probabilistic relevance framework: BM25 and beyond.
\newblock \textit{Foundations and Trends in Information Retrieval},
  3(4):333--389.

\bibitem[Sarthi et~al.(2024)]{sarthi2024raptor}
Sarthi, P., Abdullah, S., Tuli, A., Khanna, S., Goldie, A., and Manning, C.~D.
  (2024).
\newblock RAPTOR: Recursive abstractive processing for tree-organized retrieval.
\newblock In \textit{Proceedings of ICLR}.

\bibitem[Shi et~al.(2023)]{shi2023large}
Shi, F., Chen, X., Misra, K., Scales, N., Dohan, D., Chi, E., Sch{\"a}rli, N.,
  and Zhou, D. (2023).
\newblock Large language models can be easily distracted by irrelevant context.
\newblock In \textit{Proceedings of ICML}, pages 31210--31227.

\bibitem[Shuster et~al.(2021)]{shuster2021retrieval}
Shuster, K., Poff, S., Chen, M., Kiela, D., and Weston, J. (2021).
\newblock Retrieval augmentation reduces hallucination in conversation.
\newblock In \textit{Findings of EMNLP}, pages 3784--3803.

\bibitem[Su et~al.(2024)]{su2024roformer}
Su, J., Ahmed, M., Lu, Y., Pan, S., Bo, W., and Liu, Y. (2024).
\newblock RoFormer: Enhanced transformer with rotary position embedding.
\newblock \textit{Neurocomputing}, 568:127063.

\bibitem[Sviridenko(2004)]{sviridenko2004note}
Sviridenko, M. (2004).
\newblock A note on maximizing a submodular set function subject to a knapsack
  constraint.
\newblock \textit{Operations Research Letters}, 32(1):41--43.

\bibitem[Trivedi et~al.(2022)]{trivedi2022interleaving}
Trivedi, H., Balasubramanian, N., Khot, T., and Sabharwal, A. (2022).
\newblock Interleaving retrieval with chain-of-thought reasoning for
  knowledge-intensive multi-step questions.
\newblock In \textit{Proceedings of ACL}, pages 10014--10037.

\bibitem[Wang et~al.(2023a)]{wang2023survey}
Wang, L., Ma, C., Feng, X., Zhang, Z., Yang, H., Zhang, J., Chen, Z.,
  Tang, J., Chen, X., Lin, Y., et~al. (2023).
\newblock A survey on large language model based autonomous agents.
\newblock \textit{Frontiers of Computer Science}, 18(6):186345.

\bibitem[Wang et~al.(2023b)]{wang2023voyager}
Wang, G., Xie, Y., Jiang, Y., Mandlekar, A., Xiao, C., Zhu, Y., Fan, L., and
  Anandkumar, A. (2023).
\newblock Voyager: An open-ended embodied agent with large language models.
\newblock \textit{arXiv preprint arXiv:2305.16291}.

\bibitem[Xi et~al.(2023)]{xi2023rise}
Xi, Z., Chen, W., Guo, X., He, W., Ding, Y., Hong, B., Zhang, M., Wang, J.,
  Jin, S., Zhou, E., et~al. (2023).
\newblock The rise and potential of large language model based agents: A survey.
\newblock \textit{arXiv preprint arXiv:2309.07864}.

\bibitem[Xu et~al.(2023)]{xu2023recomp}
Xu, F.~F., Shi, W., and Choi, E. (2023).
\newblock RECOMP: Improving retrieval-augmented LMs with context compression
  and selective augmentation.
\newblock In \textit{Proceedings of ICLR 2024}.

\bibitem[Ye et~al.(2024)]{ye2024differential}
Ye, X., Axmed, M., Rossi, R., and Zhao, S. (2024).
\newblock Differential attention for LLMs.
\newblock \textit{arXiv preprint arXiv:2410.05258}.

\bibitem[Zaheer et~al.(2020)]{zaheer2020bigbird}
Zaheer, M., Guruganesh, G., Dubey, K.~A., Ainslie, J., Alberti, C.,
  Ontanon, S., Pham, P., Ravula, A., Wang, Q., Yang, L., and Ahmed, A. (2020).
\newblock Big bird: Transformers for longer sequences.
\newblock In \textit{Proceedings of NeurIPS}, pages 17283--17297.

\bibitem[Zhong et~al.(2024)]{zhong2024memorybank}
Zhong, W., Guo, L., Gao, Q., Ye, H., and Wang, Y. (2024).
\newblock MemoryBank: Enhancing large language models with long-term memory.
\newblock In \textit{Proceedings of AAAI}, pages 19724--19731.

\end{thebibliography}

\end{document}
